\documentclass[conference]{IEEEtran}

\usepackage{cite}
\usepackage{amsmath,amssymb,amsfonts}
\usepackage{graphicx}
\usepackage{booktabs}
\usepackage{array}
\usepackage{url}
\usepackage{xcolor}

\title{Notes on RAG }

\author{
\IEEEauthorblockN{Luis Mendez}
}

\begin{document}
\maketitle

\section{What is RAG?}
Retrieval-Augmented Generation (RAG) is a framework that combines retrieval-based methods with generative models to enhance the quality and relevance of generated content. RAG allows a model to access external knowledge sources, such as databases or documents, during the generation process, enabling it to produce more informed and contextually accurate responses.
It's a powerful approach that leverages the strengths of both retrieval and generation, making it particularly effective for tasks that require up-to-date information or domain-specific knowledge.

\end{document}